\documentclass[11pt]{article}
\usepackage[margin=1in]{geometry}
\usepackage{amsmath,amssymb,amsthm}
\usepackage[hidelinks]{hyperref}
\newtheorem{theorem}{Theorem}
\newtheorem{lemma}[theorem]{Lemma}
\newtheorem{proposition}[theorem]{Proposition}
\newtheorem{corollary}[theorem]{Corollary}
\theoremstyle{remark}
\newtheorem{remark}[theorem]{Remark}
\newtheorem{example}[theorem]{Example}
\DeclareMathOperator{\Nil}{Nil}
\DeclareMathOperator{\Per}{Per}
\title{Nilperiod rings and additive periods of power maps}
\author{}
\date{}
\begin{document}
\maketitle
\begin{abstract}
A ring is nilperiod if every nilpotent element is an additive period of some
power map. We prove that every nilperiod ring is an NI-ring, resolving a
conjecture of Burnette without finiteness, polynomial-identity, characteristic,
or identity assumptions. The proof shows that each nilpotent element generates
a nil ideal, using a square-zero substitution and descent through nil ideals.
We also prove that finite unital nilperiod rings are exactly the finite direct
products of finite local rings. Finally, we determine the additive periods of
all power maps on Corbas rings and obtain a corrected count for the cases with
nontrivial field automorphism.
\end{abstract}

\section{Introduction}
All rings are associative; an identity is assumed only when explicitly stated.
For a ring $R$, let $\Nil(R)$ denote its set of nilpotent elements. For $n\geq1$
write $\alpha_n(x)=x^n$ and
\[
 \Per(\alpha_n)=\{a\in R:(x+a)^n=x^n\text{ for every }x\in R\}.
\]
This is an additive subgroup of $R$, and substituting $x=0$ shows that every
one of its elements is nilpotent. Following Burnette \cite{Burnette}, the
ring $R$ is \emph{nilperiod} if
\[
 \Nil(R)=\bigcup_{n\geq1}\Per(\alpha_n).
\]
The exponent may depend on the nilpotent element. A ring is an \emph{NI-ring}
if its nilpotent elements form a two-sided ideal.

Burnette conjectured that every nilperiod ring is an NI-ring
\cite[Conjecture~2.6]{Burnette}, and proved the assertion for PI algebras
\cite[Theorem~2.9]{Burnette}. We prove the conjecture in its full stated
generality.

\begin{theorem}\label{thm:main}
In a nilperiod ring, every nilpotent element generates a nil two-sided ideal.
Consequently every nilperiod ring is an NI-ring.
\end{theorem}

The proof is elementary. If $a^2=0$ is an additive period of a power map,
the substitution $x=ar$ forces $ar$ to be nilpotent for every $r$.
Additive closure of nilpotents then makes the ideal generated by $a$ nil.
For an arbitrary nilpotent $a$, induction first treats the ideal generated
by $a^2$, after which the square-zero argument applies in the quotient.

The period condition also forces every idempotent to be central. This gives
an exact finite unital classification.

\begin{theorem}\label{thm:finite}
A nonzero finite unital ring is nilperiod if and only if it is a finite direct
product of finite local rings.
\end{theorem}

Here, as usual, a unital ring is local when its nonunits form a proper ideal.
The zero ring is nilperiod and can be included in Theorem~\ref{thm:finite}
as the empty direct product.
In the final section we apply this viewpoint to the twisted rings considered
in \cite[Example~2.2]{Burnette} and compute their periodic power maps exactly.

\section{Nil ideals generated by nilpotents}
An ideal is called \emph{nil} if each of its elements is nilpotent; no common
nilpotency exponent is required. To handle rings without identity, write
$R^1$ for the standard unitization $\mathbb Z\oplus R$, with multiplication
\[
 (m,r)(n,s)=(mn,ms+nr+rs).
\]
The ring $R$ embeds as a two-sided ideal in $R^1$. If $R$ already has an
identity, one may instead take $R^1=R$. The ideal $(a)$ generated by $a\in R$
is the set of finite sums of terms $uav$ with $u,v\in R^1$.
We do not assume that $R^1$ is nilperiod.

The following observation is already used in the proof of
\cite[Theorem~2.9]{Burnette}.

\begin{lemma}\label{lem:add}
The nilpotents of a nilperiod ring form an additive subgroup.
\end{lemma}
\begin{proof}
Negatives of nilpotents are nilpotent. Let $a,b$ be nilpotent. Choose $n$
with $(x+a)^n=x^n$ for every $x\in R$, and choose $m$ with $b^m=0$.
Then $(a+b)^n=b^n$, so $(a+b)^{nm}=0$.
\end{proof}

\begin{lemma}\label{lem:quotient}
If $R$ is nilperiod and $I$ is a nil ideal, then $R/I$ is nilperiod.
\end{lemma}
\begin{proof}
Suppose $a+I$ is nilpotent. Then $a^m\in I$ for some $m\geq1$.
As $I$ is nil, $(a^m)^t=0$ for some $t\geq1$. Hence $a$ is nilpotent
in $R$. A period identity for $a$ descends to the required identity for
$a+I$ in $R/I$.
\end{proof}

\begin{lemma}\label{lem:square-zero}
If $R$ is nilperiod and $a^2=0$, then $(a)$ is a nil ideal.
\end{lemma}
\begin{proof}
There is nothing to prove if $a=0$. Otherwise choose $n\geq2$ with
$(x+a)^n=x^n$ for all $x\in R$.
Let $r\in R^1$ and set $x=ar\in R$. Since $a^2=ax=0$, the
noncommutative expansion has only two possibly nonzero words:
\[
 (x+a)^n=x^n+x^{n-1}a.
\]
Indeed, any occurrence of $a$ before the final position is followed by
$x$ or $a$, and the resulting word vanishes. The period identity therefore
gives $x^{n-1}a=0$. Right multiplication by $r$ shows that
\begin{equation}\label{eq:multiple}
 (ar)^n=0\qquad(r\in R^1).
\end{equation}
For $u,v\in R^1$, it follows that
\[
 (uav)^{n+1}=u(avu)^n av=0.
\]
Thus every two-sided multiple $uav$ is nilpotent. By
Lemma~\ref{lem:add}, every finite sum of such multiples is nilpotent,
which proves the assertion.
\end{proof}

\begin{proof}[Proof of Theorem~\ref{thm:main}]
We prove by strong induction on $j\geq1$ that, in every nilperiod ring,
each element $a$ satisfying $a^j=0$ generates a nil ideal.
The cases $j=1$ and $j=2$ follow from $a=0$ and
Lemma~\ref{lem:square-zero}, respectively.
Let $j\geq3$, and assume the assertion for all smaller positive integers.
Then
\[
 (a^2)^{\lceil j/2\rceil}=0,\qquad \lceil j/2\rceil<j.
\]
The induction hypothesis makes $I=(a^2)$ a nil ideal. By
Lemma~\ref{lem:quotient}, $R/I$ is nilperiod, and the image $\bar a$ of
$a$ has square zero. Lemma~\ref{lem:square-zero} applied in $R/I$ shows
that $(\bar a)$ is nil. Since $I\subseteq(a)$, this ideal is $(a)/I$.

For $y\in(a)$, there is consequently an $m\geq1$ with $y^m\in I$.
Since $I$ is nil, $(y^m)^t=0$ for some $t\geq1$. Hence $y$ is
nilpotent, completing the induction.
Finally, Lemma~\ref{lem:add} gives additive closure of $\Nil(R)$,
and multiplication of a nilpotent $a$ by any ring element stays in the
nil ideal $(a)$. Thus $\Nil(R)$ is a two-sided ideal.
\end{proof}

\begin{corollary}
If $R$ is nilperiod, then $\Nil(R)$ is its largest nil ideal and
$R/\Nil(R)$ is reduced.
\end{corollary}
\begin{proof}
Every nil ideal is contained in $\Nil(R)$. If $a+\Nil(R)$ is nilpotent,
then $a^m\in\Nil(R)$ for some $m$, so $a$ is nilpotent and its image
is zero.
\end{proof}

\section{Idempotents and finite rings}
\begin{proposition}\label{prop:central}
Every idempotent in a nilperiod ring is central.
\end{proposition}
\begin{proof}
Let $e^2=e$ and $r\in R$. The element $a=er-ere$ satisfies
$a^2=0$, $ea=a$, and $ae=0$. Thus $e+a$ is idempotent.
A period identity for $a$, evaluated at $e$, gives
\[
 e+a=(e+a)^n=e^n=e.
\]
Hence $er=ere$. Similarly, $b=re-ere$ is square-zero and $e+b$ is
idempotent, so the same argument gives $re=ere$. Therefore $er=re$.
\end{proof}

\begin{lemma}\label{lem:idempotent-power}
Every element of a finite ring has an idempotent positive power.
\end{lemma}
\begin{proof}
For $x$ in a finite ring there exist positive integers $a,b$ such that
$x^{a+b}=x^a$. Choose $n\geq a$ divisible by $b$. Eventual periodicity
of the powers then gives $x^{2n}=x^n$.
\end{proof}

\begin{proof}[Proof of Theorem~\ref{thm:finite}]
Suppose first that $R$ is finite, unital, and nilperiod. By
Proposition~\ref{prop:central}, its idempotents are central.
Successively split the identity into orthogonal nonzero central idempotents
until each component contains no idempotents except zero and its identity.
The process terminates since $R$ is finite. It gives
\[
 R\cong R_1\times\cdots\times R_t.
\]
Each direct factor is nilperiod: a nilpotent in it lifts to a nilpotent
with zero coordinates in the other factors, and the corresponding period
identity restricts to that factor. Its nilpotents therefore form an ideal
by Theorem~\ref{thm:main}.

For $x\in R_i$, Lemma~\ref{lem:idempotent-power} gives $x^n=0$ or
$x^n=1_{R_i}$. Thus $x$ is either nilpotent or a unit. The nonunits of
$R_i$ are exactly its nilpotents, which form a proper ideal, so $R_i$ is
local.

Conversely, let $R$ be finite and local, with nonunit ideal $M$.
Its only idempotents are zero and one: if $e^2=e$, one of $e,1-e$ is a
unit, and $e(1-e)=0$ gives the conclusion.
For $x\in M$, an idempotent power remains in $M$, so it must be zero.
Thus $\Nil(R)=M$. Choose $n$ at least the maximum of the nilpotency
indices of the elements of $M$, and divisible by the exponent of the
finite group $R^\times$. Then
\[
 x^n=\begin{cases}0,&x\in M,\\1,&x\notin M.\end{cases}
\]
Translation by any element of $M$ preserves membership in $M$.
Consequently every element of $\Nil(R)$ is a period of $\alpha_n$.
For a finite direct product of finite local rings, choose one $n$
satisfying these requirements in all factors.
\end{proof}

\begin{example}
An NI-ring need not be nilperiod, even when finite. The upper triangular
$2\times2$ matrices over $\mathbb F_2$ have nilpotents
$\{0,E_{12}\}$, a square-zero ideal. But for $e=\operatorname{diag}(1,0)$
and $a=E_{12}$, the distinct matrices $e$ and $e+a$ are both idempotent.
Hence $(e+a)^n\ne e^n$ for every $n\geq1$.
\end{example}

\section{Power maps on Corbas rings}
Let $q=p^f$, where $p$ is prime and $f\geq1$, and choose $0\leq s<f$.
Let $\phi(a)=a^{p^s}$ be the corresponding automorphism of $\mathbb F_q$,
and put $g=\gcd(f,s)$. We allow $s=0$, so the identity automorphism has
$g=f$. Consider the Corbas ring
\begin{equation}\label{eq:corbas}
 R=\mathbb F_q\oplus\mathbb F_q,\qquad
 (a,b)(c,d)=(ac,ad+b\phi(c)).
\end{equation}
This is an associative unital ring, with identity $(1,0)$ and square-zero
ideal $N=\{0\}\oplus\mathbb F_q$. If $a\ne0$, the element $(a,b)$ has inverse
\[
 \bigl(a^{-1},-a^{-1}b\phi(a^{-1})\bigr).
\]
Thus $R$ is local and $\Nil(R)=N$.

The following calculation corrects the assertion in
\cite[Example~2.2 and Table~1]{Burnette} that nonidentity twists have no
periodic power maps. It also determines exactly which exponents occur.

\begin{theorem}\label{thm:corbas}
For the ring \eqref{eq:corbas}, set
\[
 M=\frac{q-1}{p^g-1}.
\]
Then
\[
 \Per(\alpha_n)=
 \begin{cases}
 N,&pM\mid n,\\
 \{0\},&pM\nmid n.
 \end{cases}
\]
There are exactly $p^g-1$ distinct periodic power maps on $R$.
\end{theorem}
\begin{proof}
Direct induction gives
\[
 (a,b)^n=\bigl(a^n,bC_n(a)\bigr),\qquad
 C_n(a)=\sum_{j=0}^{n-1}a^j\phi(a)^{n-1-j}.
\]
Every additive period must lie in $N$, by evaluation at zero. For $y\ne0$,
the element $(0,y)$ is a period exactly when $C_n(a)=0$ for all
$a\in\mathbb F_q$. At $a=1$ this forces $p\mid n$, and hence $n\geq2$.
The condition at $a=0$ then holds. If $a\ne0$ and $a=\phi(a)$, then
$C_n(a)=na^{n-1}$, which also vanishes. For $a\ne\phi(a)$ the geometric
sum gives
\[
 C_n(a)=\frac{a^n-\phi(a)^n}{a-\phi(a)}.
\]
Its vanishing is equivalent to $(a/\phi(a))^n=1$.
The homomorphism $a\mapsto a/\phi(a)$ on the cyclic group
$\mathbb F_q^\times$ has kernel $\mathbb F_{p^g}^\times$ and therefore
image of order $M$. Consequently $C_n$ vanishes everywhere if and only
if $p\mid n$ and $M\mid n$. Since $p\nmid M$, this is $pM\mid n$.
The argument also determines the full period group as stated.

For the count, the unit group of $R$ has exponent $E=p(q-1)$.
The displayed power formula shows that every unit has $E$th power one.
Elements $(a,0)$ realize order $q-1$, while $(1,b)$ with $b\ne0$ have
order $p$, so the exponent is exactly $E$.
All nonunits have square zero. Therefore, for $n,m\geq2$, the maps
$\alpha_n$ and $\alpha_m$ agree exactly when $n\equiv m\pmod E$:
sufficiency follows from the unit and nonunit cases, and necessity from
the definition of the unit-group exponent.
There are thus $E$ distinct maps with exponent at least two, and among
their exponent classes exactly
\[
 \frac{E}{pM}=p^g-1
\]
are periodic. The identity map $\alpha_1$ is distinct from these maps on
nonzero elements of $N$ and has no nonzero additive period.
\end{proof}

\begin{example}
For $q=4$ and the nontrivial automorphism $\phi(a)=a^2$, one has
$p=2$, $g=1$, $M=3$, and $E=6$. The power map $\alpha_6$ equals
$(1,0)$ on all units and zero on $N$, so its period group is $N$.
It is the unique distinct periodic power map on this noncommutative
sixteen-element ring.
\end{example}

\begin{thebibliography}{9}
\bibitem{Burnette}
C. Burnette,
\emph{On power maps over periodic rings},
Graduate J. Math. \textbf{10}, no.~1 (2025), 5--16.
\href{https://gradmath.org/wp-content/uploads/2025/07/GJM2025-Burnette.pdf}{Published article}.
Earlier preprint under the title \emph{On power maps over weakly periodic rings}:
\href{https://arxiv.org/abs/2207.14283}{arXiv:2207.14283}.
\end{thebibliography}
\end{document}
